% !PRINTED 137
\weberpage{137}
%<<CONT

setzen. Durch Multiplication mit $y^n$ und Division mit $a_n$ geht dann (21) über in
\begin{equation}
 y^n+\frac{a_{n-1}}{a_n}y^{n-1}+\cdots+\frac{a_1}{a_n}y+\frac{a_0}{a_n}=0,
\tag{23}
\end{equation}
und wir können auf die Wurzeln dieser Gleichung den vorhin ausgesprochenen Satz anwenden, indem wir $a_0=0$ annehmen, so dass eine der Wurzeln von (23) verschwindet. Wir erlangen so den Satz:

Man kann in der Gleichung (21) $a_0$ so klein annehmen, dass eine ihrer Wurzeln über alle Grenzen gross wird, während die anderen sich beliebig wenig von den Wurzeln der Gleichung $(n-1)$ten Grades
\[
 a_1x^{n-1}+a_2x^{n-2}+\cdots+a_{n-1}x+a_n=0
\]
unterscheiden.

Man drückt das auch so aus, dass \gesperrt{mit verschwindendem $a_0$ eine der Wurzeln von (21) unendlich wird.}

Wenn auch noch $a_1$ oder $a_1$ und $a_2$ u. s. f. verschwinden, so werden zwei oder drei Wurzeln unendlich.

Es rechtfertigt sich hierdurch der bisweilen gebrauchte Ausdruck, dass eine Gleichung $n$ten Grades durch Unendlichwerden einer Wurzel in eine Gleichung $(n-1)$ten Grades übergehe.

\begin{center}
\rule{0.16\textwidth}{0.4pt}
\end{center}
